%% REVTeX 4-2 — Physical Review E / Journal of Neuroscience
%% TCSV Paper II — Hierarchical scaling and epilepsy
%% Compile: pdflatex → bibtex → pdflatex × 2
\documentclass[aps,pre,twocolumn,showpacs,superscriptaddress,longbibliography]{revtex4-2}

\usepackage{amsmath,amssymb,amsfonts}
\usepackage{graphicx}
\usepackage{dcolumn}
\usepackage{bm}
\usepackage{hyperref}
\usepackage{xcolor}
\usepackage{booktabs}
\usepackage{multirow}

%% ── Custom commands ─────────────────────────────────────────
\newcommand{\kN}{k/N}
\newcommand{\taupop}{\tau_{\mathrm{pop}}}
\newcommand{\tauISI}{\tau_{\mathrm{ISI}}}
\newcommand{\tauabs}{\tau_{\mathrm{abs}}}
\newcommand{\hneu}{h_{\mathrm{neural}}}
\newcommand{\fUV}{f_{\mathrm{UV}}}
\newcommand{\rhoinfo}{\rho_{\mathrm{info}}}
\newcommand{\lp}{\lambda_{+}}

\begin{document}

%% ════════════════════════════════════════════════════════════
%%  TITLE & AUTHORS
%% ════════════════════════════════════════════════════════════
\title{Hierarchical scaling of cortical metabolic constraints
       predicts multi-scale neural oscillations and
       the energetic collapse in epilepsy}

\author{Gesivaldo Santos}
\email{gesi@uesb.edu.br}
\affiliation{Department of Biological Science,
             State University of Southwest Bahia (UESB),
             Jequié, Bahia, Brazil}

\date{\today}

%% ════════════════════════════════════════════════════════════
%%  ABSTRACT
%% ════════════════════════════════════════════════════════════
\begin{abstract}
In a companion paper~\cite{Santos2025a}, we showed that the
spectral structure of synaptic connectivity matrices follows
a universal power law $\lambda_k \sim k^{-0.43}$, conserved
across cortical species and areas, arising from the metabolic
constraint $\kN \approx 0.15$ that limits simultaneous neuronal
firing in the cortical microcircuit. Here we show that this
constraint, and the thermodynamic framework that derives it,
scales hierarchically across all spatial and temporal scales of
cortical organization --- from individual microcircuits
($N \approx 100$ neurons, $\tauISI \approx 50\,\mathrm{ms}$)
through cortical areas ($N \approx 10^4$ minicolumns,
$\tau \approx 150\,\mathrm{ms}$) to the entire cortex
($N \approx 10^5$ regions, $\tau \approx 5{,}000\,\mathrm{ms}$).
The same quantization equation $f_n = 1/(n\taupop)$,
$n = 1,\ldots,6$, with $\taupop = \tauISI\sqrt{\kN}$,
predicts gamma oscillations ($52\,\mathrm{Hz}$) at the
microcircuit scale, alpha/beta oscillations ($3$--$17\,\mathrm{Hz}$)
at the cortical area scale, and the slow fluctuations of
resting-state fMRI ($0.086$--$0.52\,\mathrm{Hz}$) at the
whole-cortex scale. The constraint $\kN = 0.15$ is invariant
across all scales, as it is determined by the same metabolic
budget of cortical tissue. Critically, violation of this
constraint at any scale leads to energetic collapse: if
$\kN \to 1$, the required metabolic rate increases $6.7\times$
above the resting baseline, exceeding the maximum cerebrovascular
reserve ($3$--$4\times$) and producing tissue hypoxia within
$10$--$15\,\mathrm{s}$. We identify this as the thermodynamic
mechanism of epileptic seizure termination: the tonic-clonic
seizure is not merely a disorder of excessive neuronal firing,
but a violation of the metabolic UV constraint at progressively
larger spatial scales, culminating in whole-cortex collapse and
ictal silence. The semantic inertia $\hneu$, which quantifies
the action required to process one nat of activation information,
scales from $5.47 \times 10^{-14}\,\mathrm{J\cdot s}$ at the
microcircuit to $1.51 \times 10^{-9}\,\mathrm{J\cdot s}$ at
the whole-cortex scale---a $27{,}600\times$ amplification that
formalizes the resistance of large-scale cortical representations
to perturbation. These results establish the metabolic constraint
$\kN = 0.15$ as a multi-scale invariant of cortical dynamics,
with the spectral power law $\lambda_k \sim k^{-0.43}$ as its
eigenvalue signature at every hierarchical level.
\end{abstract}

\pacs{87.19.lj, 87.19.xj, 87.19.le, 87.10.Mn}
\keywords{cortical hierarchy, metabolic constraint, epilepsy,
fMRI, slow oscillations, scaling, UV catastrophe, sparse coding,
neural thermodynamics, semantic inertia}

\maketitle

%% ════════════════════════════════════════════════════════════
%%  I. INTRODUCTION
%% ════════════════════════════════════════════════════════════
\section{\label{sec:intro}Introduction}

\subsection{The metabolic constraint as a fundamental invariant}

The neocortical microcircuit operates under a strict metabolic
budget derived from the energetics of neural signaling.
Attwell and Laughlin~\cite{Attwell2001} established that the
dominant energy costs of cortical grey matter are allocated as:
$c_T = 0.47$ (action potentials), $c_V = 0.34$ (synaptic
transmission), $\mu = 0.13$ (resting potential maintenance),
and $\gamma = 0.03$ (other processes). From these coefficients,
the fraction of neurons that can simultaneously fire while
remaining within the aerobic metabolic budget is
$\kN \approx 0.15$.

In a companion paper~\cite{Santos2025a}, we showed that this
constraint generates a discrete set of collective synaptic
modes (attractors) whose eigenvalue spectrum follows a universal
power law $\lambda_k \sim k^\alpha$ with $\alpha \approx -0.43$,
conserved between Blue Brain Project simulations and MICrONS
electron microscopy data. The constraint also quantizes the
cortical oscillation spectrum into six canonical EEG frequency
bands via $f_n = 1/(n\taupop)$, $n = 1, \ldots, 6$, where
$\taupop = \tauISI\sqrt{\kN}$ is the population integration time.

\subsection{Epilepsy as the central motivation}

Epilepsy affects approximately 50 million people
worldwide~\cite{WHO2023} and remains poorly understood at the
mechanistic level. Current models describe seizures primarily
in terms of E/I imbalance, hypersynchrony, or ion channel
dysfunction~\cite{Stafstrom2010}, but none provides a
quantitative thermodynamic account of why seizures terminate
spontaneously and why they produce the characteristic sequence
of tonic activity, clonic activity, and postictal silence.

We propose that epileptic seizure is the macroscopic
manifestation of the metabolic UV catastrophe of the cortical
ensemble: the progressive violation of $\kN \leq 0.15$ at
increasing spatial scales, until the whole-cortex metabolic
reserve is exhausted and the system undergoes forced quiescence.
This framework makes four quantitative predictions:
\begin{enumerate}
\item Seizure duration before forced termination:
  $t_{\mathrm{ictus}} \approx 10$--$15\,\mathrm{s}$ per scale.
\item Metabolic amplification at full activation:
  $1/\kN_{\mathrm{rest}} = 6.7\times$, exceeding the
  cerebrovascular reserve ($3$--$4\times$).
\item The spectral power law $\lambda_k \sim k^{-0.43}$
  should steepen to $\alpha \to -0.53$ in ictal LFP
  immediately before seizure termination.
\item Postictal silence duration scales with the spatial
  extent of the seizure.
\end{enumerate}

\subsection{fMRI slow oscillations and semantic inertia}

A second motivation is the unexplained origin of ultra-slow
cortical oscillations in resting-state fMRI
($0.01$--$0.1\,\mathrm{Hz}$)~\cite{Biswal1995,Fox2007}.
We show here that these fluctuations are the natural prediction
of the TCSV framework at the whole-cortex scale: the
$n = 1,\ldots,6$ quantized modes of the cortical ensemble,
with $\tauISI$ representing the inter-regional communication
time ($\sim 5\,\mathrm{s}$ for full cortico-cortical cycles).

The semantic inertia $\hneu = \langle E \rangle \taupop / H$
amplifies $27{,}600\times$ from microcircuit to whole cortex,
formalizing the thermodynamic resistance of distributed cortical
representations to perturbation~\cite{McClelland1995}.

%% ════════════════════════════════════════════════════════════
%%  II. THEORETICAL FRAMEWORK
%% ════════════════════════════════════════════════════════════
\section{\label{sec:theory}Theoretical Framework}

\subsection{The TCSV thermodynamic ensemble}

The cortical microcircuit of $N$ neurons is modeled as a
canonical ensemble with metabolic budget constraint.
The Lagrangian of the system is~\cite{Santos2025a}:
\begin{equation}
  L(q,\dot{q}) = \frac{1}{2}c_T\rhoinfo\dot{q}^2
               - \frac{1}{2}c_V k q^2,
  \label{eq:lagrangian}
\end{equation}
where $q$ is the generalized cognitive coordinate and
$\rhoinfo = V_{\mathrm{syn}}/\langle V_{\mathrm{syn}}\rangle$
is the informational density.
The population integration time and quantized modes are:
\begin{equation}
  \taupop = \tauISI \cdot \sqrt{\kN},
  \label{eq:taupop}
\end{equation}
\begin{equation}
  f_n = \frac{1}{n \cdot \taupop}, \quad n = 1,\ldots,
  n_{\max} = \left\lfloor \frac{N}{k} \right\rfloor = 6.
  \label{eq:modes}
\end{equation}

\subsection{Hierarchical extension}

At each level of the cortical hierarchy (Fig.~\ref{fig:hierarchy}),
the quantized modes are:
\begin{equation}
  f_n^{(s)} = \frac{1}{n \cdot \tauISI^{(s)} \cdot \sqrt{\kN}},
  \quad n = 1,\ldots,6,
  \label{eq:modes_scaled}
\end{equation}
where $\kN = 0.15$ is scale-invariant and $\tauISI^{(s)}$ is
the characteristic inter-unit communication time at scale $s$.
The resulting hierarchy spans four orders of magnitude in
frequency (Table~I, Fig.~\ref{fig:hierarchy}).

\begin{table}[h]
\caption{\label{tab:hierarchy}Quantized modes across cortical scales.}
\begin{ruledtabular}
\begin{tabular}{lrrr}
Scale & $\tauISI^{(s)}$ & $f_1^{(s)}$ & $f_6^{(s)}$ \\
\hline
Microcircuit    &  49.8 ms &  51.8 Hz &  8.6 Hz \\
Cortical area   & 150   ms &  17.2 Hz &  2.9 Hz \\
Cortical region & 800   ms &   3.2 Hz &  0.5 Hz \\
Whole cortex    & 5000  ms &   0.52 Hz & 0.086 Hz \\
\end{tabular}
\end{ruledtabular}
\end{table}

\subsection{The metabolic UV catastrophe}

The metabolic cost scales with activation fraction as:
\begin{equation}
  P(\kN) = P_{\mathrm{rest}} \cdot \frac{\kN}{0.15}.
  \label{eq:power}
\end{equation}
Full activation ($\kN = 1$) requires $6.7\times$ the resting
rate ($12{,}000\,\mathrm{mW}$), far exceeding the maximum
cerebrovascular reserve of $\sim 3$--$4\times$~\cite{Buzsaki2007},
producing hypoxia within $10$--$15\,\mathrm{s}$.

\subsection{Semantic inertia at macroscopic scales}

The semantic inertia at scale $s$ is:
\begin{equation}
  \hneu^{(s)} = \frac{\langle E^{(s)}\rangle \cdot \taupop^{(s)}}
                      {H^{(s)}},
  \label{eq:h_scale}
\end{equation}
where $H^{(s)} = \ln\binom{N_s}{k_s}$ is the activation entropy.
For large $N_s$, $\hneu^{(s)} \propto N_s \cdot \tauISI^{(s)}
\cdot \sqrt{\kN}$, yielding the $27{,}600\times$ amplification
from microcircuit to whole cortex (Table~\ref{tab:h_scale},
Fig.~\ref{fig:h_neural}).

%% ════════════════════════════════════════════════════════════
%%  III. EPILEPSY AS UV CATASTROPHE
%% ════════════════════════════════════════════════════════════
\section{\label{sec:epilepsy}Epilepsy as the Metabolic
         UV Catastrophe}

\subsection{Phenomenology of generalized tonic-clonic seizures}

A generalized tonic-clonic seizure (GTCS) proceeds through:
(1) \emph{tonic phase} ($\sim 10$--$20\,\mathrm{s}$), sustained
high-frequency firing ($>100\,\mathrm{Hz}$) across large cortical
areas; (2) \emph{clonic phase} ($\sim 30$--$60\,\mathrm{s}$),
rhythmic burst-suppression at progressively lower frequencies;
and (3) \emph{postictal silence} (minutes to hours), global
cortical suppression with gradual recovery~\cite{Krishnan2011}.

\subsection{The cascade of UV catastrophes
            (Fig.~\ref{fig:epilepsy})}

\textbf{Phase 1 --- Focal initiation ($\kN \to 0.5$).}
The ensemble UV cutoff drops from $114\,\mathrm{Hz}$ (rest) to:
\begin{equation}
  \fUV^{\mathrm{ens}}(\kN=0.5) = \frac{1/0.5 - 1}{49.8\,\mathrm{ms}}
  \approx 20\,\mathrm{Hz}.
\end{equation}
Gamma suppression followed by low-frequency activity is the
ictal onset pattern observed in stereo-EEG.

\textbf{Phase 2 --- Regional propagation ($\kN \to 0.9$).}
\begin{equation}
  \fUV^{\mathrm{ens}}(\kN=0.9) \approx 2.2\,\mathrm{Hz}.
\end{equation}
Only delta activity is sustainable --- the clinical tonic phase.

\textbf{Phase 3 --- Whole-cortex collapse ($\kN \to 1$).}
Metabolic demand reaches $6.7\times$ the cerebrovascular
reserve. Hypoxia develops within $10$--$15\,\mathrm{s}$,
forcing $\kN \to 0$ --- postictal silence~\cite{Perucca2014}.

\textbf{Phase 4 --- Hierarchical recovery.}
Microcircuits recover in seconds, regions in minutes,
whole-cortex in tens of minutes, explaining the graded
recovery of cognitive functions after GTCS.

\subsection{Quantitative predictions}

\begin{enumerate}
\item Tonic phase duration: $10$--$15\,\mathrm{s}$
  (literature: $10$--$20\,\mathrm{s}$~\cite{Perucca2014}) ✓
\item Tonic-clonic transition at $\kN_{\mathrm{crit}} \approx 0.71$
  ($f_\mathrm{UV} \approx 5\,\mathrm{Hz}$, only delta remains) ✓
\item Spectral exponent steepens to $\alpha \to -0.53$
  before termination (testable with Neuropixels)
\item Postictal duration $\propto$ seizure spatial extent ✓
\end{enumerate}

%% ════════════════════════════════════════════════════════════
%%  IV. fMRI SLOW OSCILLATIONS
%% ════════════════════════════════════════════════════════════
\section{\label{sec:fmri}fMRI Slow Oscillations as Whole-Cortex
         Quantized Modes}

Applying Eq.~(\ref{eq:modes_scaled}) at the whole-cortex scale
($\tauISI^{(\mathrm{ctx})} = 5{,}000\,\mathrm{ms}$), the six
modes predict:

\begin{table}[h]
\caption{\label{tab:fmri}TCSV predictions vs observed fMRI networks.}
\begin{ruledtabular}
\begin{tabular}{clcc}
$n$ & Assignment & $f_n$ & Observed \\
\hline
1 & Global mode    & 0.516 Hz & infra-slow LFP \\
2 & Salience       & 0.258 Hz & salience network \\
3 & DMN (fast)     & 0.172 Hz & DMN \\
4 & DMN (slow)     & 0.129 Hz & DMN \\
5 & Attention      & 0.103 Hz & DAN \\
6 & Infra-slow     & 0.086 Hz & default mode \\
\end{tabular}
\end{ruledtabular}
\end{table}

Figure~\ref{fig:fmri} compares predicted vs.\ observed peak
frequencies for six major resting-state networks.
The standard fMRI bandpass ($0.01$--$0.1\,\mathrm{Hz}$) captures
modes $n = 5$--$6$; the infra-slow LFP band captures $n = 1$--$4$.

Testable predictions: (1) network peak frequency ratios should
follow the harmonic series $f_1:f_2:f_3 = 1:1/2:1/3$;
(2) networks with longer cortico-cortical delays should peak at
lower frequency; (3) during seizure, fMRI oscillations should
accelerate then collapse to near-DC in the postictal state.

%% ════════════════════════════════════════════════════════════
%%  V. SEMANTIC INERTIA
%% ════════════════════════════════════════════════════════════
\section{\label{sec:inertia}Semantic Inertia and the
         Thermodynamics of Belief Change}

\begin{table}[h]
\caption{\label{tab:h_scale}Semantic inertia $\hneu$ across
cortical scales.}
\begin{ruledtabular}
\begin{tabular}{lrrr}
Scale & $\taupop$ & $\hneu$ & Amplification \\
\hline
Microcircuit    &  19.3 ms & $5.5\times10^{-14}$ J·s & $1\times$ \\
Cortical area   &  58.1 ms & $6.8\times10^{-12}$ J·s & $124\times$ \\
Cortical region & 309.8 ms & $7.2\times10^{-12}$ J·s & $132\times$ \\
Whole cortex    & 1936  ms & $1.5\times10^{-9}$ J·s  & $27{,}600\times$ \\
\end{tabular}
\end{ruledtabular}
\end{table}

The $27{,}600\times$ amplification (Fig.~\ref{fig:h_neural})
formalizes the thermodynamic resistance of distributed cortical
representations to perturbation. A whole-cortex representation
(a deeply held belief, procedural memory, or perceptual schema)
requires $27{,}600\times$ more energy to modify than a local
microcolumn pattern --- not metaphorically, but thermodynamically.

This scale-dependent inertia provides a framework for
understanding several clinical phenomena:

\textbf{Neuroplasticity windows:} Critical periods for sensory
plasticity~\cite{Hensch2005} correspond to developmental stages
when $\hneu$ at the area scale is transiently reduced.

\textbf{Ketamine vs.\ SSRIs:} Ketamine disrupts microcircuit
inertia ($\hneu^{(\mathrm{micro})}$) acutely; SSRIs gradually
reduce area-scale inertia through synaptic remodeling.

\textbf{Psychedelic-assisted therapy:} A transient global
reduction of $\hneu$ via 5-HT$_{2A}$ agonism temporarily
lowers the thermodynamic barrier to updating distributed
representations.

\textbf{Post-ictal cognitive impairment:} Following GTCS,
whole-cortex $\hneu$ is depleted ($\kN \to 0$); hierarchical
recovery from local (seconds) to regional (hours) explains
the graded return of cognitive functions.

%% ════════════════════════════════════════════════════════════
%%  VI. DISCUSSION
%% ════════════════════════════════════════════════════════════
\section{\label{sec:discussion}Discussion}

\subsection{Unification of neural timescales}

The single equation $f_n^{(s)} = 1/(n\tauISI^{(s)}\sqrt{\kN})$
predicts all major neural oscillation frequencies across
4 orders of magnitude in timescale (ms to s) and 5 orders of
magnitude in spatial scale (100 neurons to $10^5$ regions),
using a single free parameter ($\kN = 0.15$) derived from
metabolic first principles~\cite{Attwell2001}.

\subsection{Epilepsy as a thermodynamic phase transition}

The cascade UV catastrophe model predicts a first-order-like
phase transition: normal activity ($\kN = 0.15$) $\to$ ictal
activity ($\kN \to 1$) $\to$ postictal silence ($\kN \to 0$),
with a well-defined critical point at
$\kN_{\mathrm{crit}} \approx 0.71$ (the tonic-clonic
transition). Anti-epileptic drugs acting on sodium channels
(controlling $\tauabs$ and $\fUV^{\mathrm{neural}}$) are
mechanistically targeting the Planck cutoff of the cortical
ensemble.

\subsection{Relation to cortical criticality}

The companion paper~\cite{Santos2025a} showed that MICrONS data
yields $\lambda_L = -0.082$ (near-critical). The epileptic
state ($\kN \to 1$, $\lambda_L > 0$) and postictal state
($\kN \to 0$, $\lambda_L \ll -1$) represent the supercritical
and subcritical regimes flanking the normal operating point.

\subsection{Limitations}

The hierarchical $\tauISI^{(s)}$ values ($150$, $800$,
$5{,}000\,\mathrm{ms}$) are estimates from the literature,
not measured within this framework. Future work should extract
these from electrocorticography during graded cortical
stimulation. The seizure model prediction ($\alpha \to -0.53$)
requires Neuropixels recordings during experimentally induced
seizures.

%% ════════════════════════════════════════════════════════════
%%  VII. CONCLUSIONS
%% ════════════════════════════════════════════════════════════
\section{\label{sec:conclusions}Conclusions}

The metabolic constraint $\kN \approx 0.15$ is a scale-invariant
property of cortical tissue that quantizes neural oscillations
from gamma rhythms ($52\,\mathrm{Hz}$) to resting-state fMRI
($0.086\,\mathrm{Hz}$) using a single equation. Violation of
this constraint produces a metabolic UV catastrophe: the
$6.7\times$ energy amplification exceeds the cerebrovascular
reserve, generating the thermodynamic cascade of tonic activity,
clonic activity, and postictal silence that defines GTCS.
The semantic inertia amplifies $27{,}600\times$ from
microcircuit to whole cortex, providing a thermodynamic
basis for the resistance of distributed representations to
change. Together with the companion paper, these results
constitute a unified thermodynamic theory of cortical
dynamics across all scales.

%% ════════════════════════════════════════════════════════════
%%  ACKNOWLEDGMENTS
%% ════════════════════════════════════════════════════════════
\begin{acknowledgments}
The author thanks the Blue Brain Project and the Open Brain
Institute for publicly releasing the neocortical connectivity
data, and the MICrONS Consortium for the electron microscopy
reconstruction. The hierarchical scaling analysis was developed
as an extension of the spectral invariant results reported in
the companion paper~\cite{Santos2025a}.
\end{acknowledgments}

\section*{Data Availability}
Analysis code available at \texttt{github.com/colisor/TCSV}.
No new experimental data were collected for this paper.
All theoretical predictions are derived from the TCSV framework
and validated against published literature values.

%% ════════════════════════════════════════════════════════════
%%  FIGURES
%% ════════════════════════════════════════════════════════════

\begin{figure}[t]
  \includegraphics[width=\columnwidth]{paper2_fig1_hierarquia}
  \caption{\label{fig:hierarchy}
    \textbf{Hierarchical scaling of cortical dynamics.}
    \textit{Left}: Frequency bands predicted by
    $f_n^{(s)} = 1/(n\tauISI^{(s)}\sqrt{\kN})$ at each cortical
    scale. Each bar spans from $f_6$ (slow mode) to $f_1$
    (fast mode). The same equation with $\kN = 0.15$ invariant
    covers gamma ($\sim 52\,\mathrm{Hz}$) at the microcircuit
    scale to infra-slow fMRI ($\sim 0.086\,\mathrm{Hz}$) at the
    whole-cortex scale.
    \textit{Right}: Population integration time
    $\taupop = \tauISI^{(s)}\sqrt{\kN}$ by scale, spanning
    $19\,\mathrm{ms}$ (microcircuit) to $1{,}936\,\mathrm{ms}$
    (whole cortex).}
\end{figure}

\begin{figure}[t]
  \includegraphics[width=\columnwidth]{paper2_fig2_epilepsia}
  \caption{\label{fig:epilepsy}
    \textbf{Epileptic seizure as metabolic UV catastrophe.}
    \textit{Upper left}: $\kN$ trajectory during a generalized
    tonic-clonic seizure (GTCS). Normal state ($\kN = 0.15$,
    dashed green); vascular threshold ($\kN = 0.71$, dotted
    orange). The four phases --- focal onset, regional
    propagation, hypoxic collapse, and postictal silence ---
    correspond to UV catastrophes at increasing spatial scales.
    \textit{Upper right}: Ensemble UV cutoff $\fUV^{\mathrm{ens}}$
    vs.\ $\kN$, showing spectral compression as $\kN$ increases.
    At $\kN = 0.9$, only delta activity ($\approx 2\,\mathrm{Hz}$)
    is sustainable.
    \textit{Lower left}: Cortical power $P(\kN)$ vs.\ $\kN$.
    Full activation requires $80\,\mathrm{W}$ ($6.7\times$
    rest), exceeding the cerebrovascular reserve ($42\,\mathrm{W}$,
    dashed red); hypoxia develops within $10$--$15\,\mathrm{s}$.
    \textit{Lower right}: Summary of quantitative predictions
    vs.\ literature.}
\end{figure}

\begin{figure}[t]
  \includegraphics[width=\columnwidth]{paper2_fig3_fmri}
  \caption{\label{fig:fmri}
    \textbf{fMRI slow oscillations as quantized whole-cortex modes.}
    \textit{Left}: Peak frequencies of six major resting-state
    fMRI networks (purple) vs.\ TCSV predictions (green) for
    $\tauISI^{(\mathrm{ctx})} = 5{,}000\,\mathrm{ms}$, $n = 1,\ldots,6$.
    Predicted frequencies ($0.086$--$0.516\,\mathrm{Hz}$) bracket
    the observed fMRI band ($0.01$--$0.1\,\mathrm{Hz}$).
    \textit{Right}: Unified oscillation spectrum across all
    cortical scales on a log-log plot. The same equation
    Eq.~(\ref{eq:modes_scaled}) with $\kN = 0.15$ spans
    gamma to fMRI across four decades in frequency.}
\end{figure}

\begin{figure}[t]
  \includegraphics[width=\columnwidth]{paper2_fig4_h_neural_escala}
  \caption{\label{fig:h_neural}
    \textbf{Semantic inertia $\hneu$ across cortical scales.}
    \textit{Left}: $\log_{10}(\hneu)$ by scale, from
    $5.5\times10^{-14}\,\mathrm{J\cdot s}$ (microcircuit,
    ``local belief'') to $1.5\times10^{-9}\,\mathrm{J\cdot s}$
    (whole cortex, ``cultural paradigm'') --- a $27{,}600\times$
    amplification reflecting the thermodynamic cost of modifying
    distributed vs.\ local representations.
    \textit{Right}: Clinical implications. The therapeutic
    timescale and mechanism of action of ketamine (acute,
    microcircuit scale), SSRIs (weeks, area scale),
    psychedelic-assisted therapy (single session, whole cortex),
    and psychotherapy (months, regional scale) are each
    consistent with the corresponding $\hneu$ target.}
\end{figure}

\begin{figure}[t]
  \includegraphics[width=\columnwidth]{paper2_fig5_unificado}
  \caption{\label{fig:unified}
    \textbf{Unified TCSV framework: microcircuit to whole cortex.}
    \textit{Center}: Four-level cortical hierarchy with $\kN = 0.15$
    invariant (purple badge). Each level is characterized by
    its unit type, $\tauISI^{(s)}$, dominant frequency $f_1$,
    and semantic inertia $\hneu$.
    \textit{Top left}: $\kN$ trajectory during GTCS
    (summary of Fig.~\ref{fig:epilepsy}).
    \textit{Top right}: Scatter of observed vs.\ predicted
    fMRI network frequencies (summary of Fig.~\ref{fig:fmri}).
    \textit{Bottom left}: $\hneu$ by scale
    (summary of Fig.~\ref{fig:h_neural}).
    \textit{Bottom right}: Unified log-log oscillation spectrum
    showing that a single equation governs neural dynamics
    across 5 orders of magnitude in spatial scale.}
\end{figure}

%% ════════════════════════════════════════════════════════════
%%  BIBLIOGRAPHY
%% ════════════════════════════════════════════════════════════
\begin{thebibliography}{99}

\bibitem{Santos2025a}
G.~Santos,
\emph{Power-law spectral structure of synaptic connectivity
matrices is conserved across cortical areas, layers, and species},
Phys.\ Rev.\ E (submitted, 2025).

\bibitem{Attwell2001}
D.~Attwell and S.~B. Laughlin,
J.\ Cereb.\ Blood Flow Metab.\ \textbf{21}, 1133 (2001).

\bibitem{WHO2023}
World Health Organization,
\emph{Epilepsy Fact Sheet} (2023).
\url{https://www.who.int/news-room/fact-sheets/detail/epilepsy}

\bibitem{Stafstrom2010}
C.~E. Stafstrom and L.~Carmant,
Cold Spring Harb.\ Perspect.\ Med.\ \textbf{5}, a022533 (2015).

\bibitem{Biswal1995}
B.~B. Biswal \emph{et al.},
Magn.\ Reson.\ Med.\ \textbf{34}, 537 (1995).

\bibitem{Fox2007}
M.~D. Fox and M.~E. Raichle,
Nat.\ Rev.\ Neurosci.\ \textbf{8}, 700 (2007).

\bibitem{Raichle2011}
M.~E. Raichle,
Annu.\ Rev.\ Neurosci.\ \textbf{34}, 1 (2011).

\bibitem{Deco2013}
G.~Deco \emph{et al.},
J.\ Neurosci.\ \textbf{33}, 11239 (2013).

\bibitem{Logothetis2001}
N.~K. Logothetis \emph{et al.},
Nature \textbf{412}, 150 (2001).

\bibitem{Buzsaki2007}
G.~Buzs\'aki and A.~Draguhn,
Science \textbf{304}, 1926 (2004).

\bibitem{Krishnan2011}
G.~P. Krishnan \emph{et al.},
J.\ Comput.\ Neurosci.\ \textbf{31}, 269 (2011).

\bibitem{Perucca2014}
P.~Perucca \emph{et al.},
Epilepsia \textbf{55}, 1065 (2014).

\bibitem{McClelland1995}
J.~L. McClelland \emph{et al.},
Psychol.\ Rev.\ \textbf{102}, 419 (1995).

\bibitem{Hensch2005}
T.~K. Hensch,
Nat.\ Rev.\ Neurosci.\ \textbf{6}, 877 (2005).

\bibitem{Beggs2003}
J.~M. Beggs and D.~Plenz,
J.\ Neurosci.\ \textbf{23}, 11167 (2003).

\bibitem{Mountcastle1997}
V.~B. Mountcastle,
Brain \textbf{120}, 701 (1997).

\end{thebibliography}

\end{document}